\begingroup
\scriptsize
\setlength{\tabcolsep}{3pt}
\setlength\LTleft{0pt}
\setlength\LTright{0pt}
\newlength{\FixedParamTableWidth}
\setlength{\FixedParamTableWidth}{\dimexpr\textwidth-8\tabcolsep\relax}
\begin{longtable}{@{}>{\raggedright\arraybackslash}p{0.11\FixedParamTableWidth}>{\raggedright\arraybackslash}p{0.06\FixedParamTableWidth}>{\raggedright\arraybackslash}p{0.10\FixedParamTableWidth}>{\raggedright\arraybackslash}p{0.48\FixedParamTableWidth}>{\raggedright\arraybackslash}p{0.25\FixedParamTableWidth}@{}}
\caption{Fixed numerical settings used by the current O$_2$--glucose supply-demand model and calibration workflow. Values are reported on the natural scale.}
\label{tab:model_fixed_parameters}\\

\toprule
Parameter & Value & Unit & Description & Reference \\
\midrule
\endfirsthead

\toprule
Parameter & Value & Unit & Description & Reference \\
\midrule
\endhead

$\tau_{O_2}$ & $0.1$ & day & Oxygen relaxation timescale. & Fixed numerical relaxation setting; empirical workflow choice rather than a literature-derived parameter. \\

$O_{2,\mathrm{cap}}$ & $5.0$ & \% O$_2$ & Maximum oxygen level used to cap the supply-demand oxygen trajectory, inherited from the $S_0$ upper bound. & Bertout et al. (2008) \cite{bertout_impact_2008}; McKeown (2014) \cite{mckeown_defining_2014}. \\

$O_{2,\min}$ & $0$ & \% O$_2$ & Lower oxygen floor used in the logarithmic oxygen target relation. & Physical/numerical anoxia floor used to bound oxygen values; model constraint rather than literature-estimated threshold. \\

$N_{\mathrm{ref}}$ & $10^{6}$ & cells & Reference viable burden scale used in the oxygen supply-demand relation. & Oxygen-feedback normalization constant; empirical model scaling value rather than a measured biological quantity. \\

$N_{\mathrm{unit}}$ & $22$ & chromosomes & Chromosome count corresponding to ploidy 1. & Model convention defining one haploid autosome set as ploidy 1; definition rather than literature-estimated parameter. \\

$N_{\min}$ & $22$ & chromosomes & Minimum chromosome-count state retained in the discrete state space. & Discrete state-space lower bound anchored to the model ploidy convention; grid-design choice. \\

$N_{\max}$ & $154$ & chromosomes & Maximum chromosome-count state retained in the discrete state space. & Bielski et al. (2018) \cite{bielski_genome_2018}; Zack et al. (2013) \cite{zack_pan-cancer_2013}. \\

$\Delta t$ & $0.05$ & day & Simulation time step. & Simulation time step selected for numerical resolution and stability; computational setting rather than biological parameter. \\

$K$ & $10^{12}$ & cells & Carrying capacity in the logistic crowding term. & Waclaw et al. (2015) \cite{waclaw_spatial_2015}. \\

$N(0)$ & $10^{6}$ & cells & Initial total tumor population size. & Initial-condition scale for the simulation/calibration workflow; empirical setup choice. \\

$\sigma_{\mathrm{ploidy}}$ & $0.12$ & unitless & Fixed endpoint chromosome-number observation noise in the harvest likelihood. & Fixed likelihood observation-noise term; empirical measurement-error setting in the calibration workflow. \\

$\varepsilon_{\log B}$ & $10^{-12}$ & unitless & Numerical offset added before the burden log transform. & Numerical log-offset to prevent $\log(0)$; computational safeguard. \\

$N_{\min,\mathrm{pop}}$ & $10^{-12}$ & unitless & Minimum population floor used for numerical stability. & Population floor used to avoid numerical underflow or zero states; computational safeguard. \\

\bottomrule
\end{longtable}
\endgroup
